\section{Superseded proof architecture}
\label{sec:legacy-route-note}

Earlier drafts organized the arithmetic side through a frozen archimedean self
term, a critical-$J$ approximate functional equation, a one-sided
common-output interpretation of the packetized HLZ diagonal, and a compressed
Fourier-shift decomposition of the resulting direct carriers.  That
architecture is not used anywhere in the proof of
Theorem~\ref{thm:main}.

The conditional route keeps the packet polarization in the actual lag
$d=y-x$, retaining the physical additive carrier $e(-ne^d)$.  The scalar HLZ
diagonal is not asserted to remove that carrier.  Instead,
Hypothesis~\ref{hyp:one-chi-relative-defect} compares the resulting true
operator with the positive reference model only through the negative part of
their relative spectral defect. A later audit found that the claimed
holomorphic elimination of the Hermitian archimedean self block was invalid;
Lemma~\ref{lem:holomorphic-self-elimination} now records the correct
orientation identity. The conditional decomposition must account for that
block separately.

The source files implementing the older frozen-HLP/common-output route have
been removed from the active manuscript rather than reproduced here, so that
obsolete labels cannot enter the dependency graph accidentally.  They remain
recoverable from the Git history preceding the lagwise log-only repair.  This
appendix is therefore documentary only and supplies no lemma used by the main
proof.
